\documentclass[12pt]{article}

\input{../../_assets/preambulo.tex}

\begin{document}
	
	% 1. Foto de fondo
	% 2. Título
	% 3. Encabezado Izquierdo
	% 4. Color de fondo
	% 5. Coord x del titulo
	% 6. Coord y del titulo
	% 7. Fecha
	\newcommand{\R}{{\mathbb{R}}} % Reales
	\newcommand{\K}{{\mathbb{K}}} % Cuerpo arbitrario
	
	\input{../../_assets/portada}
	\portadaExamen{ffccA4.jpg}{Geometría I\\Examen XVI}{Geometría I. Examen XVI}{MidnightBlue}{-8}{28}{2026}{Arturo Olivares Martos}
	
	\begin{description}
		\item[Asignatura] Geometría I.
		\item[Curso Académico] 2020-21.
		\item[Grado] Doble Grado de Ingeniería Informática y Matemáticas.
		\item[Grupo] Único.
		%\item[Profesor] Ana María Hurtado Cortegana y Antonio Ros Mulero.
		\item[Descripción] Parcial 2.
		\item[Fecha] 14 de enero de 2021.    
	\end{description}
	\newpage

	\begin{ejercicio}
		Sea $f : V \to V'$ una aplicación lineal. Definimos
		\Func{\ol{f}}{V/\ker(f)}{\Ima(f)}{v + \ker(f)}{f(v)}
		\begin{enumerate}
			\item Demostrar que $\ol{f}$ es una aplicación (está bien definida).
			\item Demostrar que $\ol{f}$ es un isomorfismo de espacios vectoriales.
		\end{enumerate}
	\end{ejercicio}

	\begin{ejercicio}
		Sobre el espacio vectorial $\cc{M}_2(\bb{R})$ de las matrices cuadradas reales de orden 2 se definen
		\begin{gather*}
			\varphi\left(\begin{array}{cc}
				x & y \\
				z & t
			\end{array}\right) = x + t \\
			\psi\left(\begin{array}{cc}
				x & y \\
				z & t
			\end{array}\right) = x + y + z + t.
		\end{gather*}
		\begin{enumerate}
			\item Ampliar $\{\varphi, \psi\}$ a una base del espacio dual.
			\item Hallar la matriz de cambio de base entre la base dual de la usual de $\cc{M}_2(\bb{R})$ y la base obtenida en el apartado anterior.
		\end{enumerate}
	\end{ejercicio}
	
	\newpage
	\setcounter{ejercicio}{0} % Reseteo de contador para ejercicios resueltos
	
\begin{ejercicio}
	Sea $f : V \to V'$ una aplicación lineal. Definimos
	\Func{\ol{f}}{V/\ker(f)}{\Ima(f)}{v + \ker(f)}{f(v)}
	\begin{observacion}
		Notemos que se nos está pidiendo que demostremos el primer teorema de isomorfía para espacios vectoriales.
	\end{observacion}
	\begin{enumerate}
		\item Demostrar que $\ol{f}$ es una aplicación (está bien definida).
		
		Sean $x,y \in V$ y tal que $x + \ker(f) = y + \ker(f)$. Veamos que la aplicación no depende del representante elegido; es decir, que $\ol{f}(x + \ker(f)) = \ol{f}(y + \ker(f))$.
		\begin{align*}
			x + \ker(f) = y + \ker(f)
			&\iff x - y \in \ker(f)
			\iff f(x - y) = 0\iff \\ &
			\iff f(x) - f(y) = 0 
			\iff f(x) = f(y)\iff \\ &
			\iff \ol{f}(x + \ker(f)) = \ol{f}(y + \ker(f))
		\end{align*}

		Por lo tanto, $\ol{f}$ es una aplicación.
		\item Demostrar que $\ol{f}$ es un isomorfismo de espacios vectoriales.
		
		En primer lugar, hemos de ver que $\ol{f}$ es una aplicación lineal. Sean $v_1, v_2 \in V$ y $a,b \in K$, donde $K$ es el cuerpo sobre el que están definidos los espacios vectoriales. Entonces,
		\begin{align*}
			\ol{f}(a(v_1 + \ker(f)) + b(v_2 + \ker(f)))
			&= \ol{f}((av_1 + bv_2) + \ker(f)) = f(av_1 + bv_2) = \\ &
			= af(v_1) + bf(v_2) = a\ol{f}(v_1 + \ker(f)) + b\ol{f}(v_2 + \ker(f))
		\end{align*}

		Ahora, veamos que $\ol{f}$ es un monomorfismo.
		\begin{align*}
			\ker(\ol{f}) &= \{v + \ker(f) \in V/\ker(f) \mid \ol{f}(v + \ker(f)) = 0\}
			=\\&= \{v + \ker(f) \in V/\ker(f) \mid f(v) = 0\}
			=\\&= \{v + \ker(f) \in V/\ker(f) \mid v \in \ker(f)\}
			=\\&= \{v + \ker(f) \in V/\ker(f) \mid v + \ker(f) = 0 + \ker(f)\}
			=\\&= \{0 + \ker(f)\}
			= \{0\}
		\end{align*}

		Por lo tanto, $\ol{f}$ es un monomorfismo. Veamos ahora que $\ol{f}$ es un epimorfismo calculando para ello la dimensión de su imagen.
		\begin{align*}
			\dim(V/\ker(f)) &= \cancelto{0}{\dim(\ker(\ol{f}))} + \dim(\Ima(\ol{f})) 
			=\\&= \dim(V)-\dim(\ker(f)) = \dim(\Ima(f))
		\end{align*}
		Por lo tanto, $\dim(\Ima(\ol{f})) = \dim(\Ima(f))$. Además, $\Ima(\ol{f}) \subseteq \Ima(f)$, por lo que $\Ima(\ol{f}) = \Ima(f)$ y $\ol{f}$ es un epimorfismo. En consecuencia, $\ol{f}$ es un isomorfismo de espacios vectoriales.
	\end{enumerate}
\end{ejercicio}

\begin{ejercicio}
	Sobre el espacio vectorial $\cc{M}_2(\bb{R})$ de las matrices cuadradas reales de orden 2 se definen
	\begin{gather*}
		\varphi\left(\begin{array}{cc}
			x & y \\
			z & t
		\end{array}\right) = x + t \\
		\psi\left(\begin{array}{cc}
			x & y \\
			z & t
		\end{array}\right) = x + y + z + t.
	\end{gather*}
	\begin{enumerate}
		\item Ampliar $\{\varphi, \psi\}$ a una base del espacio dual.
		
		Consideramos la base usual de $\cc{M}_2(\bb{R})$, $\cc{B}_0 = \{E_{11}, E_{12}, E_{21}, E_{22}\}$, donde:
		\begin{gather*}
			E_{11} = \left(\begin{array}{cc}
				1 & 0 \\
				0 & 0
			\end{array}\right), \quad
			E_{12} = \left(\begin{array}{cc}
				0 & 1 \\
				0 & 0
			\end{array}\right), \quad
			E_{21} = \left(\begin{array}{cc}
				0 & 0 \\
				1 & 0
			\end{array}\right), \quad
			E_{22} = \left(\begin{array}{cc}
				0 & 0 \\
				0 & 1
			\end{array}\right)
		\end{gather*}

		Sea ahora la base dual de la base usual, $\cc{B}_0^* = \{\varphi_{11}, \phi_{12}, \phi_{21}, \phi_{22}\}$, donde:
		\begin{gather*}
			\varphi_{ij}\begin{pmatrix}
				a_{11} & a_{12} \\
				a_{21} & a_{22}
			\end{pmatrix} = a_{ij}\qquad i,j \in \{1,2\}
		\end{gather*}

		En ese sistema de coordenadas, tenemos:
		\begin{gather*}
			\varphi \equiv (1,0,0,1)_{\cc{B}_0^*} \\
			\psi \equiv (1,1,1,1)_{\cc{B}_0^*}
		\end{gather*}
		
		Veamos ahora que $\{\varphi, \psi, \varphi_{21}, \varphi_{22}\}$ forman una base de $(\cc{M}_2(\bb{R}))^*$. \begin{align*}
			\begin{vmatrix}
				1 & 1 & 0 & 0 \\
				0 & 1 & 0 & 0 \\
				0 & 1 & 1 & 0 \\
				1 & 1 & 0 & 1
			\end{vmatrix} \neq 0
		\end{align*}
		Por lo tanto, $\cc{B}^*=\{\varphi, \psi, \varphi_{21}, \varphi_{22}\}$ es una base de $(\cc{M}_2(\bb{R}))^*$.
		\item Hallar la matriz de cambio de base entre la base dual de la usual de $\cc{M}_2(\bb{R})$ y la base obtenida en el apartado anterior.
		
		Tenemos que:
		\begin{align*}
			M_{\cc{B}^*\leftarrow \cc{B}_0^*} = \left(M_{\cc{B}_0^*\leftarrow \cc{B}^*}\right)^{-1}
			= \begin{pmatrix}
				1 & 1 & 0 & 0\\
				0 & 1 & 0 & 0\\
				1 & 1 & 1 & 0\\
				1 & 1 & 0 & 1
			\end{pmatrix}^{-1}
		\end{align*}

		Para calcular la inversa, empleamos el método de Gauss-Jordan.
		\begin{multline*}
			\left(\begin{array}{cccc|cccc}
				1 & 1 & 0 & 0 & 1 & 0 & 0 & 0\\
				0 & 1 & 0 & 0 & 0 & 1 & 0 & 0\\
				0 & 1 & 1 & 0 & 0 & 0 & 1 & 0\\
				1 & 1 & 0 & 1 & 0 & 0 & 0 & 1
			\end{array}\right)
			\xrightarrow{F'_4 = F_4-F_1}
			\left(\begin{array}{cccc|cccc}
				1 & 1 & 0 & 0 & 1 & 0 & 0 & 0\\
				0 & 1 & 0 & 0 & 0 & 1 & 0 & 0\\
				0 & 1 & 1 & 0 & 0 & 0 & 1 & 0\\
				0 & 0 & 0 & 1 & -1 & 0 & 0 & 1
			\end{array}\right)
			\\ \xrightarrow[F'_3 = F_3-F_2]{F'_1 = F_1-F_2}
			\left(\begin{array}{cccc|cccc}
				1 & 0 & 0 & 0 & 1 & -1 & 0 & 0\\
				0 & 1 & 0 & 0 & 0 & 1 & 0 & 0\\
				0 & 0 & 1 & 0 & 0 & -1 & 1 & 0\\
				0 & 0 & 0 & 1 & -1 & 0 & 0 & 1
			\end{array}\right)
		\end{multline*}

		Por tanto, tenemos que:
		\begin{align*}
			M_{\cc{B}^*\leftarrow \cc{B}_0^*} = \left(M_{\cc{B}_0^*\leftarrow \cc{B}^*}\right)^{-1}
			= \begin{pmatrix}
				1 & 1 & 0 & 0\\
				0 & 1 & 0 & 0\\
				1 & 1 & 1 & 0\\
				1 & 1 & 0 & 1
			\end{pmatrix}^{-1}
			= \begin{pmatrix}
				1 & -1 & 0 & 0\\
				0 & 1 & 0 & 0\\
				0 & -1 & 1 & 0\\
				-1 & 0 & 0 & 1
			\end{pmatrix}
		\end{align*}
	\end{enumerate}
\end{ejercicio}

\end{document}
	